\subsection{Model Architectures}
The evaluated models are grouped by their input processing and output
representation. This distinguishes temporal scalar regressors from frame-based
and spatial predictors, while treating physics-informed losses and uncertainty
estimation as separate methodological components.

\subsubsection{Temporal CNN--LSTM Architectures}

\paragraph{CNNLSTM (Custom Lightweight)}
The custom CNNLSTM applies three convolutional blocks with 16, 32, and 64
channels to each frame. Adaptive average pooling produces a feature vector for
each frame, and an LSTM with hidden size 64 aggregates the frame features over
time. A fully connected regression head maps the last LSTM state to the four
sensor temperatures.

\paragraph{Pretrained CNNLSTM}
This model uses an ImageNet-pretrained ResNet-18~\cite{he2016deep} as a frame encoder, followed
by an LSTM with hidden size 128 and a fully connected regression head. For the
five-channel input, the first convolution is adapted: its RGB weights are
initialized from the pretrained layer and the additional optical-flow channel
weights are initialized separately. The regression head uses the last LSTM
state to predict the four sensor temperatures.

\paragraph{PhysicsCNNLSTM}
PhysicsCNNLSTM also uses frame-wise CNN feature extraction followed by an LSTM,
but predicts sensor values at every time step rather than only from the last
state. Its encoder is either a pretrained ResNet-18 or a compact CNN, and its
regression head maps each LSTM output to four sensor values. The physical
regularization used during training is described separately below; it does
not change the network's CNN--LSTM structure.

\subsubsection{Frame-Based and Spatial ResNet Predictors}

\paragraph{SimpleResNet}
SimpleResNet uses a ResNet-18~\cite{he2016deep} encoder with an adapted first convolution for
the input channel count and a regression head for the four sensor
temperatures. It has no temporal state: when given a sequence-shaped input,
the forward pass selects the last frame.

\paragraph{SpatialResNet}
SpatialResNet retains the spatial feature maps of a ResNet-18~\cite{he2016deep} by removing its
classification and global-pooling layers. A convolutional decoder and
bilinear upsampling produce a temperature map. For a sequence input, each
frame is processed independently and the maps are returned as a sequence;
there is no recurrent temporal aggregation in this model.

\subsubsection{Spatiotemporal Spatial Architecture}

\paragraph{ConvectionBioheat}
The model evaluated under the name \emph{ConvectionBioheat} is implemented by
SpatialPhysicsCNNLSTM. A ResNet-18~\cite{he2016deep} encoder extracts spatial features from
each frame. The flattened feature maps are processed by an LSTM, whose output
is mapped to a coarse $4\times4$ temperature field at each time step. A
separate parameter head, evaluated from the first frame's features, predicts
$4\times4$ maps for the effective coefficients $\alpha$ and $\beta$. The
network returns the temperature field together with these two parameter
maps. The Bioheat residual and its optional optical-flow term are part of the
training loss, not additional network architectures.

\subsubsection{KAN-Based Hybrid Architectures}
\label{sec:kan_models}
The KAN models combine a convolutional ResNet-18 encoder with a
Kolmogorov--Arnold Network (KAN) head. KANs replace fixed node activations
with learnable univariate functions on network edges~\cite{liu2024kan}.
KAN-based approaches have also been explored for physics-informed learning
and spatiotemporal fields~\cite{wang2024kinn,patra2024pikan,jacob2025spikan};
recurrent encoders followed by KAN heads provide a related hybrid
design~\cite{shen2025alpkan}. In the B-spline formulation used here, an edge
function is represented by a weighted sum of spline basis functions:
\begin{align}
    y_j &= \sum_{i=1}^{d_{in}}
    \left( w^{b}_{ji}\,\sigma(x_i) +
    w^{s}_{ji}\,\phi_{ji}(x_i) \right), \\
    \phi_{ji}(x_i) &= \sum_{k=1}^{G+p} c_{jik}\,B_{k,p}(x_i).
\end{align}
Here, $\sigma$ is the SiLU base activation, $B_{k,p}$ are B-spline basis
functions, and the coefficients and weights are learned. We use the
efficient implementation described in~\cite{blealtan2024efficientkan}.

\paragraph{KAN-ResNet}
KAN-ResNet replaces the scalar regression head of the ResNet-18 encoder with
a two-layer KAN ($512\rightarrow64\rightarrow4$), producing the four sensor
temperatures. For sequence inputs, the encoded frame features are averaged
over time before entering the KAN head.

\paragraph{Spatial KAN-Bioheat}
Spatial KAN-Bioheat uses a separable KAN field head, following the
SPIKAN-style factorization of coordinate functions~\cite{jacob2025spikan}, to produce a coarse
spatiotemporal temperature field. Separate one-dimensional KANs evaluate
spatial and temporal coordinates; their coordinate features are combined
with the ResNet encoder representation by a fusion KAN. The resulting
$4\times4$ field at each time step is trained with the same Advanced BioHeat
Loss as ConvectionBioheat, providing a comparison of the two field
parameterizations under a shared physics-informed objective.

\subsubsection{Convolutional Liquid Time-Constant Architecture}
\label{sec:conv_ltc}
Liquid Time-Constant networks provide a recurrent model with input-dependent
state dynamics~\cite{lechner2020neural}. The evaluated LTC model, ConvLTC,
consists of a two-layer convolutional
encoder, a convolutional Liquid Time-Constant (LTC) recurrent core, and a
convolutional decoder. The encoder maps each frame to feature maps while
preserving their spatial dimensions. The recurrent state likewise retains
the spatial layout.

At each time step, the recurrent cell concatenates the current encoded frame
with the previous hidden state. Convolutions generate an input-dependent
leak and drive signal, and an Euler update advances the hidden state. The
decoder maps each recurrent state to a temperature map, producing a
spatiotemporal sequence of predictions. In the current implementation, the
update uses a fixed unit time step; irregular sampling intervals are not
passed to the recurrent core.


\subsection{Physics-Informed Training Objectives}

\subsubsection{Temporal Regularization for PhysicsCNNLSTM}
PhysicsCNNLSTM is trained with a temporal residual based on Newton's law of
cooling:
\begin{equation}
    \frac{dT}{dt} = -k(T - T_{env}).
\end{equation}
The implementation approximates the time derivative using frame-to-frame
differences and combines the residual penalty with the data loss:
\begin{equation}
    \mathcal{L}_{total} =
    \mathcal{L}_{MSE} + \lambda\mathcal{L}_{physics},\qquad
    \mathcal{L}_{physics} =
    \frac{1}{N}\sum_{i=1}^{N}
    \left|\frac{dT}{dt}+k(T-T_{env})\right|^2.
\end{equation}
Because the measurements are acquired during active MWA heating, this
cooling-law residual omits the microwave heat source and is treated as a
phenomenological temporal regularizer, not as a mechanistic model of the
heating process.

\subsubsection{Advanced BioHeat Loss}
The Advanced BioHeat Loss is applied to spatial predictions. For a
spatiotemporal temperature field, its reduced residual combines the temporal
change, spatial diffusion, a perfusion-like relaxation term, and an optional
constant source term:
\begin{equation}
    \frac{\partial T}{\partial t}
    + \vec{v}\cdot\nabla T
    = \beta\nabla^2 T
    - \alpha(T-T_{arterial}) + Q_{met}.
\end{equation}
The advection-like term is included only when optical flow $\vec{v}$ is
provided. Optical flow from B-mode video is an image-motion surrogate, not a
direct measurement of tissue velocity or fluid convection. The constant
$Q_{met}$ is learnable when enabled; it is not separately identifiable as
metabolic heating or microwave absorption. The experiments use avascular agar
phantoms, and the effective coefficients are not calibrated physical
measurements.

For single-frame spatial predictions, no temporal derivative can be
computed, so the loss uses a spatial residual instead. SpatialPhysicsCNNLSTM
also returns $\alpha$ and $\beta$ maps from its parameter head; these are
passed to the loss as learned effective maps. These terms describe
configurations of the objective and model outputs, rather than separate
architectures.

\subsection{Uncertainty Quantification}
Uncertainty is estimated using Deep Ensembles and Bayesian Neural Networks.
These methods modify the estimation procedure or regression layers and are
described separately from the deterministic backbone families.

\subsubsection{Deep Ensembles}
We train $M=5$ independent instances of SimpleResNet with different random
seeds. The ensemble prediction is the mean of the individual outputs; the
between-model variance estimates epistemic uncertainty. This approach
requires multiple models at inference time.

\subsubsection{Bayesian Neural Networks}
Bayesian layers represent weights by probability distributions and are
trained using variational inference, following the Bayes-by-Backprop
formulation~\cite{blundell2015weight}. The objective combines the data loss
with a Kullback--Leibler penalty:
\begin{equation}
    \mathcal{L}_{BNN} =
    \mathcal{L}_{MSE}
    + \beta\,\mathrm{KL}\!\left(q_\theta(w)\,\|\,P(w)\right).
\end{equation}
The KL term regularizes the approximate posterior $q_\theta(w)$ toward the
prior $P(w)$. The repository includes a Bayesian-head ResNet, a full Bayesian
ResNet variant, and a Bayesian CNNLSTM. In the Bayesian-head model, the
ResNet-18 feature extractor is deterministic and the regression head uses
Bayesian linear layers; the full Bayesian variant also replaces convolutional
layers. BayesianCNNLSTM applies its Bayesian head to each LSTM time step.
Multiple stochastic forward passes are used to estimate predictive means and
variances.

\paragraph{Bayesian Physics-Informed Combination}
Bayesian models can also be trained with a physics-informed objective. In
this case, the data term, physics residual, and KL regularizer are combined:
\begin{equation}
    \mathcal{L}_{total} =
    \mathcal{L}_{MSE}
    + \lambda_{phys}\mathcal{L}_{bioheat}
    + \lambda_{KL}\mathrm{KL}\!\left(q_\theta(w)\,\|\,P(w)\right).
\end{equation}
This combination does not by itself guarantee calibrated uncertainty or
physical consistency; those properties must be assessed empirically.
